\documentclass[10pt,a4paper]{amsart}
\usepackage[margin=19mm]{geometry}
\usepackage{amsmath,amssymb,mathtools}
\usepackage[scr=rsfs,cal=euler]{mathalfa}
\usepackage{xfrac}
\usepackage[unicode,hidelinks,bookmarks=false]{hyperref}
\newcommand{\ie}{i.e.\ }
\newcommand{\cf}{cf.\ }
\newcommand{\resp}{respectively\ }
\newcommand{\Subsection}{\subsection}
\providecommand{\nobreakdash}{\nobreak\hbox{-}}
\setlength{\emergencystretch}{2em}
\multlinegap0pt
\usepackage{amssymb}
\usepackage[shortlabels]{enumitem}
\newtheorem{theorem}{Theorem}
\title{Hyperbolicity and GCD for $n+1$ divisors with non-empty intersection}
\author{Julie Tzu-Yueh Wang, Zheng Xiao}
\date{Source: arXiv:2506.03534v2 (CC BY 4.0)}
\begin{document}
\maketitle

\noindent\emph{Source and licence.} Hyperbolicity and GCD for $n+1$ divisors with non-empty intersection. Version arXiv:2506.03534v2 (CC BY 4.0), distributed by arXiv under the Creative Commons Attribution 4.0 International licence, http://creativecommons.org/licenses/by/4.0/, as recorded in the arXiv licence metadata for this exact version and shown on the abstract page https://arxiv.org/abs/2506.03534v2. Full licence notice: \texttt{licenses/arxiv-2506.03534-CC-BY-4.0.txt}.\par\smallskip

\noindent\textit{Editorial scope.} Counted unit: the article's theorem complexgcd2 (source0.tex lines 663--697). The article's complete original proof of it is delivered with it (source0.tex lines 726--850, one \texttt{proof} environment of 125 lines, which the article places away from the statement). No mathematics is written, repaired or re-derived: the statement and the proof are the article's own lines, in the article's own notation, and no retained line is substituted or re-flowed.  Retained uncounted as the source-local context the proof needs: lines 271--285.  Nothing was omitted from any retained span, so the package carries no omission record.  No retained line contains a citation, so the wrapper carries no bibliography and no bibliography entry.  No retained line contains a figure environment, an image-inclusion command or drawing code, so no image is delivered and the package declares no asset dependency.  Editorial formatting only, in addition: an AMS \texttt{amsart} preamble that restates the article's own declarations and macros used by the retained lines, namely the environments \texttt{theorem} on separate counters, together with the standard packages those lines need.  The retained environments print in this document as Theorem 1 in order of appearance, whereas the article prints the same items with the numbering of its own full document.  The complete native source of the article as distributed in the arXiv e-print for this exact version is bundled unchanged as \texttt{sources/179272/source0.tex}.
 \begin{theorem}\label{trungeneral}
    Let $X$ be a complex projective variety of dimension $n$, and let $D_1,\cdots,D_q$  be effective Cartier divisors on  $X$.
  Let $\mathcal M$ be a finite collection of pairs $(I,Y)$, where 
$I \subset \{1,\ldots,q\}$ has cardinality $n-1$, $Y$ is a closed subscheme of $\bigcap_{i \in I} D_i$, and the divisors $D_i$, $i \in I$, intersect properly  such that $ \bigcap_{i \in I} D_i\supset Y $
is a  regular sequence. Let $\mathscr L$ be a big line sheaf on $X$.  
 Then, for every $\varepsilon>0$,  there exists a proper  Zariski-closed subset $Z\subset X$ such that for any non-constant holomorphic map $f:\mathbb{C}\to X$ with  $f(\mathbb{C})\not\subset Z$,  the following  inequality  holds:
 \begin{align}\label{mainineq}
\int_0^{2\pi} \max_{(I,Y) \in \mathcal{M}} \Bigg(
&\sum_{i \in I} \beta(\mathscr{L}, D_i)\lambda_{D_i}(f(re^{i\theta}))  
 + \Big( \beta(\mathscr{L}, Y) - \sum_{i \in I} \beta(\mathscr{L}, D_i) \Big) \lambda_Y(f(re^{i\theta}))
\Bigg)   \frac{d\theta}{2\pi}\notag\\
&\leq_{\operatorname{exc}} (1+\varepsilon) T_{\mathscr{L}, f}(r),
\end{align}
where the maximum is taken over all   $(I,Y)\in\mathcal M$, and the notation $\leq_{\operatorname{exc}}$ means the inequality holds for all $r\in \mathbb{R}^+$ except a set of finite Lebesgue measure.
\end{theorem}

\begin{theorem}\label{complexgcd2}
Let $X$ be a complex projective variety  of dimension $n$. 
Let $D_1, \ldots,D_{n+1}$ be effective Cartier  divisors on $X$ such that  there exist positive integers $a_1,\ldots, a_{n+1}$ for which $a_1D_1, \ldots, a_{n+1}D_{n+1}$ are all numerically equivalent to an ample divisor $D$.  
Suppose that  any collection of $n$ divisors among  $D_1, \ldots,D_{n+1}$ intersect properly, and that 
\begin{align*}
\bigcap_{i=1}^{n+1}D_i\neq \emptyset.
\end{align*}
Define
\begin{align*}
\beta_0=\min_{\substack{|I|=n\\Q\in \bigcap_{i=1}^{n+1}D_i }} \beta(D,(\bigcap_{i \in I} a_iD_i)_Q),
\end{align*}
where $I$ range over all possible subsets of $\{1,\ldots,n+1\}$ with $|I|=n$.
Furthermore, suppose that for every point $Q\in \bigcap_{i=1}^{n+1}D_i$ and any two subsets $I,I'$ of the indices $\{1,\ldots,n+1\}$ with $|I|=|I'|=n$,  the following inequality holds:
\begin{align}
\label{beta1cond}
(\beta(D,(\bigcap_{i \in I}a_iD_i)_Q)-1)+(\beta(D,(\bigcap_{i \in I'}a_iD_i)_Q)-1)((n+1)\beta_0-n)>0. 
\end{align}
In particular,  note that 
 $\beta_0\geq \frac{n}{n+1}$ and condition \eqref{beta1cond} is satisfied if the stronger assumption
\begin{align}
\label{beta2cond}
\beta(D,(\bigcap_{i \in I}a_iD_i)_Q)>1
\end{align}
%or
%\begin{align}
%\label{49ineq2}
%(D_i)^I_Q<(\frac{n}{n+1})^n(D_i)^I
%\end{align}
holds for all $Q\in  \bigcap_{i=1}^{n+1}D_i  $ and all $I \subset \{1,\ldots,n+1\}, |I|=n$.  
For each $\varepsilon>0$,   there exists a proper Zariski-closed subset $Z\subset X$ such that for any non-constant holomorphic map $f:\mathbb{C}\to X\setminus{\bigcup_{i=1}^{n+1} } D_i$ with $f(\mathbb{C})\not\subset Z$,  we have 
\begin{align*}
T_{f, \,\bigcap_{i=1}^{n+1}a_iD_i}(r) 
        \le \varepsilon  T_{f,D}(r)+O(1). 
\end{align*}
\end{theorem}

 \begin{proof}[Proof of Theorem \ref{complexgcd2}]
The strategy of this proof is to take a convex combination of two applications of Theorem \ref{trungeneral}, taking different sets of $n$ divisors to cut  out schemes with a common point in their support.

For simplicity of notation, we will again write $D_i=a_iD_i$ throughout the proof.    Under this convention, $D_i$'s are numerically equivalent to $D$. Given a point  $Q\in \bigcap_{i=1}^{n+1}D_i $, $z\in\mathbb C$ and  a subset $I \subset \{1,\ldots,n+1\}$ with $|I|=n$,  we introduce the following shorthand notation, which will be used throughout the proof
:\begin{align}\label{shortnotation}
\beta_{I,Q}=\beta(D,(\bigcap_{i \in I} D_i)_{Q}),\quad\text{and}\quad
\lambda_{I,Q}(f(z))&=\lambda_{(\bigcap_{i \in I} D_i)_{Q}}(f(z)). 
\end{align}


Since $D_i\equiv D$, we have $\beta(D,D_i)=\frac{1}{n+1}$ for all $i$, and it follows that for all $Q\in X$ and any $I \subset \{1,\ldots,n+1\}$ with $|I|=n$, we have
\begin{align*}
\beta_{I,Q}\geq \beta(D,\bigcap_{i \in I}D_{i})\geq \sum_{i \in I}\beta(D,D_i)\geq \frac{n}{n+1}.
\end{align*}
In particular, $\beta_0\geq \frac{n}{n+1}$.
Now set 
$$
b = \frac{1}{(n+1)\beta_0-n+1}.
$$ Then $0<b\leq 1$.   By the definition of $\beta_0$, for any $I \subset \{1,\ldots,n+1\}$ with $|I|=n$, and $Q\in  \bigcap_{i=1}^{n+1} D_{i}  $, 
\begin{align}\label{bbeta}
b ( (n+1)\beta_{I,Q}  -n+1 ) \geq 1. 
\end{align}
Define
\begin{align*}
\gamma=(n+1)\min_{\substack{I\neq J,|I|=|J|=n\\Q\in  \bigcap_{i=1}^{n+1} D_{i}  }} b(\beta_{I,Q}-1)+(1-b)(\beta_{J,Q} -1).
\end{align*}
By hypothesis \eqref{beta1cond} and the fact that $b>0$, it follows that $\gamma>0$.  Fix a point $Q\in \bigcap_{i=1}^{n+1}D_i $, and $z\in\mathbb C$.  For any $I \subset \{1,\ldots,n+1\}$ with $|I|=n$, define
\begin{align*}
\beta_{I} :=\beta_{I,Q}, \quad\text{and}\quad
\lambda_{I} :=\lambda_{I,Q}(f(z)). 
\end{align*}
Assume there exist two such index sets  $I,J$, and suppose that $\lambda_{I}\geq \lambda_{J}$.  We then compute
\begin{align}\label{gammabeta}
&b((n+1)\beta_{I}-n+1)\lambda_{I}+(1-b)((n+1)\beta_{J}-n+1)\lambda_{J}\cr
&\geq \lambda_{I}+\left(b((n+1)\beta_{I}-n+1)-1+(1-b)((n+1)\beta_{J}-n+1) \right)\lambda_{J}\cr
&= \lambda_{I}+\left(1+  (n+1)\big(b(\beta_{I}-1)+(1-b)(\beta_{J}-1)  \big) \right) \lambda_{J}\cr
&\geq \lambda_{I}+(1+\gamma)\lambda_{J}.
\end{align}

Let $P=f(z)\in X$.  Choose a permutation
\[
\{1_z,\ldots,(n+1)_z\}=\{1,\ldots,n+1\}
\]
(depending on $z$) such that
\[
\lambda_{D_{1_z}}(f(z)) \ge \lambda_{D_{2_z}}(f(z)) \ge \cdots \ge \lambda_{D_{(n+1)_z}}(f(z)).
\]
Define the two index sets
\[
I_z:=\{1_z,\ldots,n_z\},\qquad 
J_z:=\{1_z,\ldots,(n-1)_z,(n+1)_z\}.
\]
For each $z\in \mathbb C$, there exists a point $Q_z\in \text{Supp} (\bigcap_{i \in I_z} D_{i})$ (depending on $f(z)$) such that
\begin{align*}
\lambda_{\bigcap_{i \in I_z} D_{i}}(f(z))=\lambda_{ (\bigcap_{i \in I_z} D_{i})_{Q_z}}(f(z))+O(1),
\end{align*}
where the constant in the $O(1)$ is independent of $f(z)$.
We claim 
\begin{align}\label{localWeilQz0}
 &\sum_{i=1}^{n+1}\lambda_{ D_{i}}(f(z))+ \gamma \lambda_{\bigcap_{i=1}^{n+1}D_{i}} (f(z)) \cr
& = \sum_{i=1}^{n-1}\lambda_{D_{i_z}}(f(z))  +\lambda_{   I_z,Q_z }(f(z))+  (1+\gamma) \lambda_{ J_z,Q_z }(f(z))+O(1). 
\end{align}
 

If $Q_z\not\in \text{Supp} \bigcap_{i=1}^{n+1}D_{i}$, then
\begin{align*}
\lambda_{ D_{(n+1)_z}   }(f(z))&=\min \{\lambda_{D_{(n+1)_z} }(f(z)),\lambda_{\bigcap_{i  \in I_z} D_i}(f(z))\}+O(1)\\
&=\min \{\lambda_{D_{(n+1)_z} }(f(z)),\lambda_{   I_z,Q_z }(f(z))\}+O(1)\\
&=\lambda_{(\bigcap_{i \in I_z} D_{i})_{Q_z}\cap D_{(n+1)_z}}(f(z))=O(1),  
\end{align*}
since $({\bigcap_{i \in I_z} D_{i}})_{Q_z}\cap D_{(n+1)_z}$ is empty.  
Then
\begin{align*}
  &\sum_{i=1}^{n+1}\lambda_{ D_{i}}(f(z)) + \gamma \lambda_{\bigcap_{i=1}^{n+1}D_{i}} (f(z)) 
 = \sum_{i=1}^{n}\lambda_{ D_{i_z}}(f(z)) +O(1)\cr
&= \sum_{i=1}^{n-1}\lambda_{D_{i_z}}(f(z))  +\lambda_{   I_z,Q_z }(f(z))+O(1).
\end{align*}
As $\lambda_{J_z,Q_z }(f(z))=O(1)$ in this case, the claim \eqref{localWeilQz0} follows.
When $Q_z\in \text{Supp} (\bigcap_{i=1}^{n+1} D_{i})$, %let $I_{j,z}=\{1_v,\ldots,(j-1)_z,(j+1)_z,\ldots,(n+1)_z\}$, 
we use the estimate
\begin{align*}
\lambda_{ D_{(n+1)_z} } (f(z)) &=\lambda_{(\bigcap_{i=1}^{n+1} D_{i})_{Q_z}} (f(z)) +O(1)\\
&\leq \lambda_{ J_z,Q_z }(f(z)) (f(z) )+O(1) 
\end{align*}
to conclude the calim  \eqref{localWeilQz0}.
% Then
%\begin{align}\label{localWeilQz2}
 %&\sum_{i=1}^{n+1}\lambda_{ D_{i}}(f(z)) \cr
%&\le \bigg(\sum_{i=1}^{n-1}\lambda_{D_{i_z}}(f(z)) \bigg) +\lambda_{ (\bigcap_{i \in I_z} D_{i})_{Q_z}}(f(z))+ \lambda_{(\bigcap_{i \in J_z} D_{i})_{Q_z} } (f(z) )+O(1). 
%\end{align}

Since $\lambda_{   I_z,Q_z }(f(z))\ge \lambda_{ J_z,Q_z }(f(z))+O(1)$,
we may apply \eqref{gammabeta} to \eqref{localWeilQz0} to obtain
\begin{align}\label{localWeilQz3}
 & \sum_{i=1}^{n+1}\lambda_{ D_{i}}(f(z))+ \gamma \lambda_{\bigcap_{i=1}^{n+1}D_{i}} (f(z))  \notag\\
 &\le \bigg(\sum_{i=1}^{n-1}\lambda_{D_{i_z}}(f(z)) \bigg)+b\left((n+1)\beta_{I_z,Q_z}-n+1)\lambda_{I_z,Q_z}(f(z)\right) \notag\\
&\quad+(1-b)\left((n+1)\beta_{J_z,Q_z}-n+1\right)\lambda_{J_z,Q_z}(f(z))+O(1). 
\end{align}
Let $\varepsilon>0$.  We will apply Theorem \ref{trungeneral} twice.
First, for each point $z\in\mathbb C$ we take $I=\{1_z,\ldots,(n-1)_z \}$ and $Y_z=(\bigcap_{i \in I_z} D_{i})_{Q_z} $.  Apply the theorem with this choice and multiply the resulting inequality by $b$.
Second, again for each point $z\in\mathbb C$ we take  $I=\{1_z,\ldots,(n-1)_z \}$ but now let $Y_z=(\bigcap_{i \in J_z} D_{i})_{Q_z} $   and multiply the resulting inequality by $1-b$.
Combining these two applications with \eqref{localWeilQz3}, we obtain that there exists a proper Zariski-closed subset  $Z\subset X$, independent of $f$, such that  
\begin{align}\label{proxi2}
  \sum_{i=1}^{n+1}m_f(D_i,r)  +\gamma\cdot m_f\bigg( {\bigcap_{i=1}^{n+1} }  D_{i} ,r\bigg)   
    \leq_{\operatorname{exc}} (n+1+\varepsilon)T_{D,f}(r) 
\end{align}
if $f(\mathbb{C})\not\subset Z$. 
In addition, if the image of $f$ does not intersects any $D_i$, $1\le i\le n+1$, then  $N_f(D_i,r) =O(1)$ for  $1\le i\le n+1$, and hence $N_f({\bigcap_{i=1}^{n+1} }  D_{i},r) =O(1)$.
 Note that 
\begin{equation*}
    N_f(D_i,r)+ m_f(D_i,r) =T_{D_i,f}(r)  
\ge \left(1-\frac{\epsilon}{n+1}\right) T_{D,f}(r)+O(1),
\end{equation*}
since the divisors $D_i$ are numerically equivalent to $D$. Therefore, \eqref{proxi2} yields
\begin{align}\label{ht}
(n+1-\varepsilon)T_{D,f}(r)   
+\gamma\, T_{\bigcap_{i=1}^{n+1}D_{i},f}(r)   
\leq_{\operatorname{exc}} (n+1+\varepsilon)T_{D,f}(r).
\end{align}
 Consequently,
\begin{align*} 
 T_{{\bigcap_{i=1}^{n+1} }  D_{i},f}(r)   
    \leq_{\operatorname{exc}} \frac{2\varepsilon}{\gamma} T_{D,f}(r).
\end{align*}
\end{proof}

\end{document}
